% paper/sections/theorem2_lower_bound.tex -- P2-03, plan/03-phase2-theory.md.
%
% Standalone section fragment, \input{} by paper/main.tex once P6-01 builds
% the skeleton. Same preamble requirements as theorem1_bound.tex. Depends
% on paper/sections/setup.tex and theorem1_bound.tex's notation.
%
% Part A (the change-of-measure lower bound) is a mechanical adaptation of
% Kaufmann, Cappe & Garivier (2016)'s transportation-lemma technique --
% cited from established knowledge of that paper's general structure, not
% from a copy on hand; if a reviewer's copy disagrees with this section on
% an exact constant, trust the reviewer's copy over this draft. Part B (the
% impossibility construction) is marked \needshuman per the plan's own
% explicit instruction: "Part B's construction is where the real
% mathematical content is."

\section{Theorem 2: the lower bound and the impossibility phase transition}
\label{sec:theorem2}

\subsection{Part A: a change-of-measure lower bound}
\label{sec:theorem2-parta}

\begin{theorem}[Compute-weighted characteristic time lower bound]
\label{thm:lower-bound}
For any $\delta$-correct policy (Definition~\ref{def:delta-correct}, $\delta < 1/2$) and
any instance $\nu$ (a full specification of every $\theta_k$ and $h_k \in \mathcal{H}$),
\[
  \mathbb{E}_\nu[C(\tau)] \;\ge\; \mathrm{kl}(\delta, 1-\delta) \cdot T^\star(\nu),
\]
where $\mathrm{kl}(x, y) = x \log(x/y) + (1-x)\log((1-x)/(1-y))$ and $T^\star(\nu)$
solves the compute-weighted optimal-design program
\[
  \frac{1}{T^\star(\nu)} \;=\; \sup_{w} \ \min_{k \ne k^\star} \
  \frac{\Delta_k^2}{2 \, J(\theta_k, s^\star)^\top I_k(w)^{-1} J(\theta_k, s^\star)},
  \qquad
  I_k(w) = \sum_{s \in \mathcal{S}} w(k, s) \, \frac{J(\theta_k, s) J(\theta_k, s)^\top}{\sigma^2(s)},
\]
the supremum over allocations $w$ (a measure on (arm, scale) pairs with
$\sum_{k,s} w(k,s)\, c(s) = 1$). This program is concave in $w$ (Lemma~\ref{lem:concave}
below) and therefore both a well-posed optimization problem and numerically solvable --
exactly the program P3-02 solves.
\end{theorem}

\begin{remark}[Where this departs from classical best-arm identification]
\label{rem:not-classical-bai}
In classical fixed-confidence BAI (Kaufmann, Capp\'e \& Garivier, 2016; Garivier \&
Kaufmann, 2016), each pull of arm $k$ reveals a sample with a per-arm
Kullback-Leibler divergence to any alternative instance, and the characteristic time
is $1/T^\star = \sup_w \min_{k} w_k \, d(\mu_k, \mu_{k^\star})$-shaped (a per-arm KL
rate). Here, no scale $s < s^\star$ ever directly observes $\mu_k(s^\star)$: information
about it arrives only through the parametric link $g(\theta_k, \cdot)$, so a pull's value
is governed by how much it shrinks the \emph{Fisher information matrix} $I_k(w)$ around
$\theta_k$, propagated to $s^\star$ through the Jacobian -- the
$(J^\top I_k(w)^{-1} J)^{-1}$-shaped rate in place of a per-arm KL is the structural
novelty relative to the classical setting, and is exactly what makes scale
\emph{allocation} (not just arm allocation) the object P3-02 needs to optimize.
\end{remark}

\subsubsection{Proof}

\paragraph{Step 1 -- the transportation lemma (cited, standard).}
For any two instances $\nu, \nu'$ and any event $E$ measurable with respect to the
policy's observations, the standard change-of-measure / data-processing inequality for
sequential experiments (Kaufmann, Capp\'e \& Garivier, 2016, Lemma 1, following the
classical Wald-type argument) gives
\[
  \sum_{k,s} \mathbb{E}_\nu[N_k(s, \tau)] \, \mathrm{KL}\big(P_{k,s}^\nu, P_{k,s}^{\nu'}\big)
  \;\ge\; \mathrm{kl}\big(\mathbb{P}_\nu[E], \mathbb{P}_{\nu'}[E]\big),
\]
where $N_k(s,\tau)$ counts pulls of $(k,s)$ up to the stopping time and $P_{k,s}^\nu$ is
the observation distribution at $(k,s)$ under $\nu$. \emph{Cited, not re-derived}: this
is foundational BAI machinery, established well before this project and not specific to
the parametric setting here.

\paragraph{Step 2 -- applying it to a confusing alternative.}
Fix $k \ne k^\star$. Take $E = \{\hat{k} = k^\star\}$, so $\delta$-correctness gives
$\mathbb{P}_\nu[E] \ge 1-\delta$. Construct $\nu'$ identical to $\nu$ except
$\theta'_k = \theta_k + \Delta\theta$ for a perturbation chosen so that $k$ is optimal
under $\nu'$ (i.e. $g(\theta'_k, s^\star) > \mu_{k^\star}(s^\star)$, at least
marginally); $\delta$-correctness under $\nu'$ then gives
$\mathbb{P}_{\nu'}[E] = \mathbb{P}_{\nu'}[\hat{k}=k^\star] \le \delta$. Plugging into
Step 1's inequality: $\mathrm{kl}(\mathbb{P}_\nu[E], \mathbb{P}_{\nu'}[E]) \ge
\mathrm{kl}(1-\delta, \delta) = \mathrm{kl}(\delta, 1-\delta)$ ($\mathrm{kl}$ is
symmetric in this sense). Mechanical, given Step 1.

\paragraph{Step 3 -- the per-pull KL, linearized.}
For Gaussian (or, to the same local-quadratic order, general sub-Gaussian)
observations at $(k,s)$ with mean $g(\theta_k, s)$ and variance $\sigma^2(s)$,
$\mathrm{KL}(P_{k,s}^\nu, P_{k,s}^{\nu'}) = (g(\theta_k,s) - g(\theta'_k,s))^2 /
(2\sigma^2(s)) \approx (\Delta\theta^\top J(\theta_k,s))^2 / (2\sigma^2(s))$, exact for
linear-in-$\theta$ $g$ and a first-order (in $\Delta\theta$) approximation otherwise --
the same nonlinear-$g$ caveat as Theorem~\ref{thm:main-bound} Step 1
(\needshuman{the exact same non-asymptotic gap: this expansion is exact as
$\Delta\theta \to 0$ / $C \to \infty$ for a general nonlinear family, matching this
theorem's role as an \emph{asymptotic} ($\delta \to 0$) characteristic-time result
(Theorem~\ref{thm:algorithm-correctness}'s own "asymptotic optimality" framing is
exactly this), but a fully explicit non-asymptotic version needs the same M-estimator
technology flagged there, not a new derivation}). Summing over the design:
$\sum_s \mathbb{E}_\nu[N_k(s,\tau)] \, \mathrm{KL}(\cdot) \approx \mathbb{E}_\nu[C(\tau)]
\cdot \Delta\theta^\top I_k(w) \Delta\theta / 2$ where $w(k,s) = \mathbb{E}_\nu[N_k(s,\tau)]
\, c(s) / \mathbb{E}_\nu[C(\tau)]$ is the (compute-normalized) empirical allocation.
Mechanical algebra, given Step 3's linearization.

\paragraph{Step 4 -- optimizing the confusing alternative, and over $k$.}
Combining Steps 2-3: $\mathbb{E}_\nu[C(\tau)] \ge \mathrm{kl}(\delta,1-\delta) \big/
\big(\Delta\theta^\top I_k(w) \Delta\theta / 2\big)$ for \emph{any} $\Delta\theta$ with
$J(\theta_k,s^\star)^\top \Delta\theta \ge \Delta_k$ (the minimum shift needed to flip
the gap at $s^\star$). The adversary (seeking the \emph{weakest} such bound, i.e. the
\emph{cheapest} confusing alternative) minimizes $\Delta\theta^\top I_k(w) \Delta\theta$
subject to this constraint; by the standard quadratic-form identity
$\min\{x^\top A x : b^\top x = c\} = c^2 / (b^\top A^{-1} b)$ (Lagrange multipliers; a
D-optimal-design identity, mechanical given $I_k(w) \succ 0$), the minimum is
$\Delta_k^2 / (J^\top I_k(w)^{-1} J)$, giving exactly the per-$k$ bound in the theorem
statement. Taking the worst (smallest) bound over $k \ne k^\star$ -- since the policy
must simultaneously guard against every alternative -- and the best allocation $w$ the
policy could have used gives $T^\star(\nu)$ as stated. Mechanical.
$\blacksquare$

\begin{lemma}[Concavity of the characteristic-time program]
\label{lem:concave}
$w \mapsto \min_{k \ne k^\star} \Delta_k^2 / (2 J_k^\top I_k(w)^{-1} J_k)$ is concave.
\end{lemma}
\begin{proof}
For fixed $k$, $I_k(w)$ is linear in $w$ (Theorem~\ref{thm:lower-bound}'s definition),
so, by Step 4's identity read in reverse,
$\Delta_k^2 / (2J_k^\top I_k(w)^{-1} J_k) = \min_{\Delta\theta : J_k^\top \Delta\theta =
\Delta_k} \Delta\theta^\top I_k(w) \Delta\theta / 2$, an infimum over a family of
functions each \emph{linear} (hence concave) in $w$ -- an infimum of concave functions
is concave. The overall objective is then a minimum over $k$ of concave functions,
again concave. Mechanical (the same two-line argument used throughout the BAI lower-bound
literature for the classical, non-parametric $T^\star$ program; adapted here rather than
re-derived from scratch).
\end{proof}

\subsection{Part B: the impossibility phase transition}
\label{sec:theorem2-partb}

\begin{theorem}[Impossibility below the observed range]
\label{thm:impossibility}
Suppose $\mathcal{H}$ and $\mathcal{S}$ admit two instances $\nu_1, \nu_2$ that (i)
induce \emph{identical} observation distributions at every $(k, s)$ with $s < s^\star$,
and (ii) have $k^\star(\nu_1) \ne k^\star(\nu_2)$. Then no policy using only scales
$s < s^\star$ is $\delta$-correct simultaneously on $\nu_1$ and $\nu_2$, for any
$\delta < 1/2$, at any finite compute.
\end{theorem}

\begin{proof}
Since every $(k,s)$ observation with $s<s^\star$ has an identical distribution under
$\nu_1$ and $\nu_2$, and the policy (its choice of pulls, its stopping time $\tau$, and
its recommendation $\hat{k}$) is a measurable function of those observations alone, the
\emph{joint law} of $(\tau, \hat{k})$ is identical under $\nu_1$ and $\nu_2$. Write
$k_1^\star = k^\star(\nu_1) \ne k_2^\star = k^\star(\nu_2)$. Then
$\mathbb{P}_{\nu_2}[\hat{k} = k_1^\star] = \mathbb{P}_{\nu_1}[\hat{k} = k_1^\star] \ge
1-\delta$ (identical law; $\delta$-correctness on $\nu_1$), while also
$\mathbb{P}_{\nu_2}[\hat{k} = k_2^\star] \ge 1-\delta$ ($\delta$-correctness on $\nu_2$).
Since $k_1^\star \ne k_2^\star$, these are disjoint events under a single measure
$\mathbb{P}_{\nu_2}$, so their probabilities sum to at most $1$:
$(1-\delta) + (1-\delta) \le 1 \iff \delta \ge 1/2$. Contradiction for $\delta < 1/2$.
Mechanical -- this is the same two-point indistinguishability argument (Le Cam's method)
used in Theorem~\ref{thm:identifiability}'s minimax lower bound, applied here to a
\emph{decision} rather than an estimation error.
\end{proof}

\begin{remark}[The genuinely hard part: does such a pair exist?]
\label{rem:impossibility-construction}
Theorem~\ref{thm:impossibility}'s proof is entirely mechanical \emph{given} $\nu_1,
\nu_2$. \needshuman{Whether -- and precisely when -- $\mathcal{H}$ and $\mathcal{S}$
admit such a pair is the real mathematical content the source document flags as this
paper's main point of departure from classical BAI ("just BAI applied" is the
anticipated reviewer objection). What follows is a fully worked \emph{candidate}
construction for two natural choices of $\mathcal{H}$, with numerical verification that
it satisfies (i)-(ii) above, plus a partial (sufficient-condition-only) structural
characterization. A tight, general, necessary-and-sufficient characterization of the
phase transition across arbitrary $\mathcal{H}$ is not attempted here and needs a
human's judgement on how far to generalize it for the paper.}
\end{remark}

\subsubsection{A worked candidate construction}
\label{sec:theorem2-construction}

Fix the largest fitting scale $s_{\max} = \max S_{\mathrm{fit}} < s^\star$ and gap
$g = s^\star - s_{\max} > 0$. Let $\nu_1$ have $h_k \equiv 0$ (arm $k$ exactly
in-family) and $\nu_2$ have $h_k(s) = \eta_k \cdot \phi\big((s - s_{\max})/g\big)$ for a
bump function $\phi: \mathbb{R} \to [0,1]$ with $\phi(x) = 0$ for $x \le 0$ and
$\phi(1) = 1$ (so $h_k \equiv 0$ on $S_{\mathrm{fit}}$ -- both instances agree exactly
on every fitting scale -- and $h_k(s^\star) = \eta_k$). Both instances share the same
$\theta_k$'s, so the two arms' orderings differ only through $h$'s effect at $s^\star$:
choosing $\eta_k$'s that reverse the sign of $\mu_{k^\star}(s^\star) - \mu_k(s^\star)$
for the true winner $k^\star$ (under $\nu_1$) gives $k^\star(\nu_2) \ne k^\star(\nu_1)$.

\begin{itemize}
  \item \textbf{Sup-norm class} ($\mathcal{H} = \{h : \sup_s |h(s)| \le \eta\}$, no
    smoothness requirement beyond boundedness): any $\phi$ works, including a step
    function ($\phi(x) = \mathbb{1}[x > 0]$) if $\mathcal{H}$ does not require
    continuity, or any continuous bump otherwise. \textbf{Sufficient condition for
    impossibility: $\eta \ge \Delta_{\min} / 2$} where $\Delta_{\min} = \min_{k \ne
    k^\star} \Delta_k$ under $\nu_1$ (enough perturbation budget, split between the
    two competing arms, to close a gap of $\Delta_{\min}$ and reverse it). This
    condition does not depend on the gap $g$ at all -- with no smoothness tying
    behaviour at $s^\star$ to behaviour at $s_{\max}$, impossibility is purely a
    question of how much perturbation budget $\mathcal{H}$ allows, however close
    $s^\star$ is to the fitted region.
  \item \textbf{H\"older-$\alpha$ ball} ($\mathcal{H} = \{h : |h(s) - h(s')| \le
    \eta\, |s-s'|^\alpha\}$): a bump function respecting this modulus can rise from $0$
    at $s_{\max}$ to at most $\eta \, g^\alpha$ by $s^\star$ (the H\"older condition
    applied directly to the pair $(s_{\max}, s^\star)$, or, for a bump with the
    standard $\mathrm{dist}(x, \{0,1\})^\alpha$-type profile, exactly this order).
    \textbf{Sufficient condition for impossibility: $\eta \, g^\alpha \ge
    \Delta_{\min}/2$} -- unlike the sup-norm case, this is a genuine \emph{phase
    transition in $g$}: for fixed $\eta, \alpha, \Delta_{\min}$, impossibility
    (via \emph{this construction}) holds once the target is far enough past the
    fitted region ($g$ large enough), and this specific construction fails for smaller
    $g$. Whether the regime is then actually \emph{solvable} for small $g$, or merely
    that this particular construction does not witness impossibility there, is exactly
    the open direction flagged in Remark~\ref{rem:impossibility-construction}.
\end{itemize}

\subsubsection{Numerical certificate}
\label{sec:theorem2-certificate}

\texttt{tests/theory/test\_theorem2.py} builds both constructions explicitly (a concrete
$\phi$ satisfying each modulus) on a randomized grid of $(S_{\mathrm{fit}}, s^\star,
\theta_k, \eta)$ instances and checks, for each: (a) $h_1 \equiv 0$ and $h_2$ agree with
$\max_{s \in S_{\mathrm{fit}}} |h_2(s)| = 0$ exactly (machine precision) on every fitting
scale, so \emph{any} function of the fitting data alone -- in particular every one of
the six fitters in \texttt{src/pdt/scaling/fitters.py} -- produces an identical fit
under $\nu_1$ and $\nu_2$, mechanically (not something that needs a "large class of
estimators" checked one by one: the fitted $\hat\theta$ is a deterministic function of
data that is identical under both instances, full stop); (b) the true winner flips
($k^\star(\nu_1) \ne k^\star(\nu_2)$); (c) the H\"older construction's sufficient
condition ($\eta g^\alpha \ge \Delta_{\min}/2$) correctly predicts, on both sides, which
randomly generated instances succeed at flipping the winner while retaining the
Hölder-$\alpha$ modulus and which do not. Verifying $T^\star$ against the P3 algorithm's
achieved sample complexity in the solvable regime is deferred to P3-04's simulation
study, which did not exist when this task was written -- flagged here rather than
skipped silently.
